\subsection{Study area: the Rogue River and the Wolverine source}
\label{sec:study}
We demonstrate the integrated model on the Rogue River watershed (Michigan; USGS gauge~04118500),
a $606$~km\textsuperscript{2} glacial-drift tributary of the Grand River gauged at Rockford. The
watershed contains the Wolverine World Wide tannery legacy and the House Street disposal site, one
of the most intensively investigated PFAS contaminated-site corridors in the United States, with a
dense ambient surface-water network and --- decisively for this work --- a dense groundwater PFAS
plume dataset. The same SWAT\textsuperscript{+} model used for the surface-water engine
(Section~\ref{sec:engine}; $17{,}771$ HRUs, $590$ channels, generated from NHDPlus High Resolution
hydrography) supplies the land surface, recharge, and channel network; the MODFLOW~6 model uses a
uniform 250~m grid with three layers ($\approx 24{,}900$ active cells), summarized with the rest of
the coupled-model configuration in Table~\ref{tab:studyarea}. The contribution is the
\emph{automated, reproducible} construction of the coupled model and its contaminant extension on the
SWATGenX platform \citep{Rafiei2026SWATGenX}, whose high-resolution SWAT\textsuperscript{+} models run
on an accelerated, byte-reproducible engine \citep{rafiei2026engine}; the same code runs for any
U.S.\ gauge.

\begin{figure}[ht]
  \centering
  \includegraphics[width=0.62\textwidth]{fig1_study_area.pdf}
  \caption{The Rogue River watershed (USGS~04118500), Michigan. The SWATGenX-generated
  \citep{Rafiei2026SWATGenX} SWAT\textsuperscript{+} network resolves $590$ channels from NHDPlus High
  Resolution; the red star
  marks the Wolverine World Wide / House Street legacy PFAS source in the lower river. The same
  watershed carries the dense surface-water and groundwater PFAS observation networks used to
  validate both compartments. Map lines delineate study areas and do not necessarily depict accepted
  national boundaries.}
  \label{fig:study}
\end{figure}

\subsection{On-the-fly MODFLOW~6 model generation}
\label{sec:generation}

The complete groundwater model --- land surface, layered aquifer geometry, hydraulic properties,
recharge, and boundary conditions --- is generated from public data at model-build time. Building
each component on the fly, rather than from pre-computed raster products, keeps every input auditable
and guarantees consistency with the surface-water model. The land surface is the SWAT\textsuperscript{+}
DEM; aquifer geometry and conductivity are kriged from a raw well database; and recharge is taken
directly from the SWAT\textsuperscript{+} water balance. A uniform-thickness bedrock buffer beneath
the drift and water table is required for steady-state Newton stability on dissected basins, where a
sliver bedrock layer in deep-water-table uplands otherwise dewaters and oscillates.

\paragraph{Land surface.} The model top is the SWAT\textsuperscript{+} DEM resampled onto the
MODFLOW grid centres, the single source of truth for elevation with all other layers snapped beneath
it. Sharing the surface-water model's own elevation eliminates a class of elevation errors that arise
when a separate DEM product is clipped and cleaned independently and removes an external dependency.

\paragraph{Aquifer geometry and hydraulic properties.} Layer bottoms and hydraulic properties are
kriged on the fly from a raw groundwater-well database (in the demonstration region, the Michigan
Wellogic database, $\approx 1.17$~million wells; a companion effort harmonizes such well records into a
national inventory of subsurface lithology and hydraulics \citep{rafiei2026gwinventory}). Each record is cured (sentinel and quality flags
stripped, physical consistency enforced, units converted), then a variogram is fit and ordinary
kriging with $k$-fold cross-validation produces each surface: static water level, horizontal and
vertical hydraulic conductivity, and aquifer thickness. The three-layer column is snapped to the DEM
top --- an upper and lower drift layer whose combined thickness follows the kriged well-derived
aquifer bottom, and a bedrock buffer beneath the kriged water table --- so the saturated zone is
contained within the active domain everywhere. Storing the well database as a spatially indexed
FlatGeobuf reduces a bounding-box query from minutes to under a second.

\paragraph{Recharge.} Recharge is taken from the SWAT\textsuperscript{+} model itself --- the
per-HRU percolation below the soil profile --- rather than from a static recharge raster, so it
already integrates the weather, soils, land use, and evapotranspiration physics of the calibrated
surface-water model. Percolation is mapped to MODFLOW cells with the area-weighted HRU-to-cell map
described in Section~\ref{sec:coupling} and applied as a long-term annual mean (and, in the coupled
mode, as a daily series). A multiplier on recharge is retained as a calibration parameter.

\paragraph{Boundary conditions.} The watershed perimeter is a head-dependent general-head boundary
(GHB) toward the kriged water table, with conductance scaled by hydraulic conductivity and saturated
thickness. This is more defensible than a rigid constant head: it draws near-boundary heads toward
the regional water table without prescribing them and without importing the full kriging error of
the static-water-level surface. The perimeter is taken directly from the active-domain edge (every
active cell with a non-active four-neighbour or on the grid edge, in every layer), which guarantees
a complete, leak-free ring rather than the open perimeter cells and interior fixed heads that a
buffered-boundary rasterization can leave behind. Lakes are represented as a surface-seepage drain
at the land surface. The flow solver uses the Newton formulation with a backtracking,
delta-bar-delta under-relaxation configuration tuned for robust convergence at 250~m resolution.

\paragraph{Streams.} The stream network --- $1{,}506$ streamflow-routing (SFR) reaches built from
the SWAT\textsuperscript{+} channel graph, detailed in Section~\ref{sec:coupling} --- represents
only channels of stream order $\geq 2$. Order-1 flowlines are largely artificial headwater traces
that sit high in the landscape, dry cells when incised, and do not sustain regional
groundwater--stream exchange.
